\begin{figure}[!htbp]
    \centering
    \begin{tikzpicture}[
        level 1/.style={sibling distance=6cm, level distance=1.6cm},
        level 2/.style={sibling distance=3.2cm, level distance=1.6cm},
        rootnode/.style={draw, rounded corners=4pt, minimum height=0.9cm, minimum width=3.2cm, align=center, font=\small\bfseries, fill=black!8},
        archnode/.style={draw, rounded corners=4pt, minimum height=0.9cm, minimum width=3cm, align=center, font=\small},
        leafnode/.style={draw, rounded corners=3pt, minimum height=0.8cm, minimum width=2.6cm, align=center, font=\small},
        secref/.style={font=\scriptsize, text=black!50},
        edge from parent/.style={draw, ->, >=stealth, thick},
    ]
    \node[rootnode] {语音到语音生成}
        child {
            node[archnode] {编码器-解码器架构\\{\scriptsize (\S\ref{sec:rw_encdec})}}
            child {
                node[leafnode] {语音到文本翻译\\{\scriptsize (\S\ref{sec:rw_encdec_s2tt})}}
            }
            child {
                node[leafnode] {语音到语音翻译\\{\scriptsize (\S\ref{sec:rw_encdec_s2st})}}
            }
        }
        child {
            node[archnode] {仅解码器架构\\{\scriptsize (\S\ref{sec:rw_deconly})}}
            child {
                node[leafnode] {语音到文本生成\\{\scriptsize (\S\ref{sec:rw_deconly_s2t})}}
            }
            child {
                node[leafnode] {语音到语音生成\\{\scriptsize (\S\ref{sec:rw_deconly_s2s})}}
            }
        };
    \end{tikzpicture}
    \bicaption{\enspace 本章技术框架：按模型架构和任务类型组织的研究现状}{\enspace Technical framework of this chapter: research landscape organized by model architecture and task type}
    \label{fig:rw_framework}
\end{figure}
